\unit{第149次报告——凝聚代数层的对偶定理}
\sid{149}
\src{169}
\seminarzh{（1957年5月）}
\paperhead{凝聚代数层的对偶定理}%
  {亚历山大·格罗滕迪克 著}

下述结果受塞尔的\fquote{代数对偶定理}启发，发现于1955年冬和1956年冬。借助
关于射影空间上同调的若干相当初等的结果[3]，并大量采用[2]所述形式的
Cartan--Eilenberg同调代数，便可极其简明地证明这些结果。

\fsection{1}{模层的 Ext（[2]，第3、4章）}

设 $X$ 是一个拓扑空间，并赋有一个带幺环（不必交换）的层 $\Osh$。考虑由
$\Osh$-模层构成的阿贝尔范畴 $\Ccat^{\Osh}$；这些对象也称为 $\Osh$-模。已知
该范畴的每个对象都具有内射消解，因而可在其中定义具有熟知形式性质的
$\operatorname{Ext}$ 函子。为避免混淆，更确切地说，我们以
$\operatorname{Hom}_{\Osh}(X;A,B)$，或简记为 $\operatorname{Hom}(X;A,B)$，
表示从 $A$ 到 $B$ 的 $\Osh$-同态所成的阿贝尔群；而以
$\Homsh_{\Osh}(A,B)$ 表示从 $A$ 到 $B$ 的同态芽所组成的层
（$A,B\in\Ccat^{\Osh}$）。固定 $A\in\Ccat^{\Osh}$，定义函子 $h_A$ 和
$\hsh_A$；它们分别取值于阿贝尔群范畴 $\Ccat$ 和 $X$ 上阿贝尔群层的范畴
$\Ccat^{\Zsh}$，其定义式为：
\begin{equation}\tag{1.1}
  h_A(B)=\operatorname{Hom}_{\Osh}(X;A,B),
  \qquad
  \hsh_A(B)=\Homsh_{\Osh}(A,B).
\end{equation}

函子 $h_A$ 和 $\hsh_A$ 都是协变左正合函子。考察它们的右导出函子，分别记为
$\operatorname{Ext}_{\Osh}^{p}(X;A,B)$ 和 $\Extsh_{\Osh}^{p}(A,B)$。因此按定义有
\begin{equation}\tag{1.2}
\left\{
\begin{aligned}
\operatorname{Ext}_{\Osh}^{p}(X;A,B)
  &= (R^{p}h_A)(B)
   = H^{p}\bigl(\operatorname{Hom}_{\Osh}(X;A,C(B))\bigr),\\
\Extsh_{\Osh}^{p}(A,B)
  &= (R^{p}\hsh_A)(B)
   = H^{p}\bigl(\Homsh_{\Osh}(A,C(B))\bigr),
\end{aligned}
\right.
\end{equation}

其中符号 $R^{p}$ 表示取右导出函子，$C(B)$ 表示 $B$ 在 $\Ccat^{\Osh}$ 中的
任意内射消解。以
\[
  \Gamma:\Ccat^{\Zsh}\longrightarrow\Ccat
\]
表示\fquote{截面}函子。回顾其右导出函子记为
\[
  B\longmapsto H^{p}(X,B).
\]

\src{170}
\begin{equation}\tag{1.3}
  H^{p}(X,B)=(R^{p}\Gamma)(B)=H^{p}\bigl(\Gamma(C(B))\bigr).
\end{equation}
显然 $h_A=\Gamma\hsh_A$；另一方面，容易验证 $\hsh_A$ 将内射对象变为
$\Gamma$-无环对象。由熟知的结论可得：

\result{命题}{1} 对每个 $\Osh$-模 $A$，在 $\Ccat^{\Osh}$ 上存在一个上同调
谱函子，它收敛于分次函子
$\bigl(\operatorname{Ext}_{\Osh}^{*}(X;A,B)\bigr)$，且其初始项为
\begin{equation}\tag{1.4}
  E_{2}^{p,q}(A,B)=H^{p}\bigl(X,\Extsh_{\Osh}^{q}(A,B)\bigr).
\end{equation}
由此得到边缘同态和一条五项正合列，兹不写出。

\result{推论}{1} 若 $A$ 局部同构于 $\Osh^{n}$，则有典范同构
\begin{equation}\tag{1.5}
  \operatorname{Ext}_{\Osh}^{p}(X;A,B)
  \overset{\sim}{\longleftarrow}
  H^{p}\bigl(X,\Homsh_{\Osh}(A,B)\bigr)
\end{equation}
（由该谱序列的边缘同态给出）。特别地，有典范同构
\begin{equation}\tag{1.6}
  \operatorname{Ext}_{\Osh}^{p}(X;\Osh,B)=H^{p}(X,B).
\end{equation}

为应用这些结果，必须能够明确描述 $\Extsh_{\Osh}^{p}(A,B)$。这些函子可以
局部计算：若 $U$ 是 $X$ 的开子集，则
\[
  \Extsh_{\Osh}^{p}(A,B)|U
  =\Extsh_{\Osh|U}^{p}(A|U,B|U),
\]
这是因为内射 $\Osh$-模限制到 $U$ 上后仍是内射 $\Osh|U$-模。此外，固定
$x\in X$，有函子性同态
\begin{equation}\tag{1.7}
  \Homsh_{\Osh}(A,B)_x
  \longrightarrow
  \operatorname{Hom}_{\Osh_x}(A_x,B_x),
\end{equation}
它唯一延拓为关于 $B$ 的上同调函子之间的一个同态：
\begin{equation}\tag{1.8}
  \Extsh_{\Osh}^{p}(A,B)_x
  \longrightarrow
  \operatorname{Ext}_{\Osh_x}^{p}(A_x,B_x).
\end{equation}

\result{命题}{2} 若在 $x$ 的某个邻域内，$A$ 同构于某个同态
$\Osh^{m}\longrightarrow\Osh^{n}$ 的余核，则当 $A$ 局部具有由有限秩自由模组成、延伸至次数 $p+1$ 的部分消解时，(1.8) 在次数 $p$ 是同构；若 $A$ 伪凝聚，则它在每个次数都是同构。仅有有限表示只能保证零次的比较。%
\footnote{印刷文本作“(1.7)”；$p$ 的出现及前一句均要求此处为“(1.8)”。}
特别地，若 $A$ 是凝聚 $\Osh$-模 [3]，上述结论成立。

\src{171}
\result{命题}{3} 设 $L_{*}=(L_i)$ 是 $\Osh$-模 $A$ 的一个左消解，其中每个
$\Osh$-模都局部同构于某个 $\Osh^{n}$。则 $\Extsh_{\Osh}^{*}(A,B)$ 可与
$H^{*}\bigl(\Homsh_{\Osh}(L_{*},B)\bigr)$ 等同，而
$\operatorname{Ext}_{\Osh}^{*}(X;A,B)$ 可与复形
$\Homsh_{\Osh}(L_{*},B)$ 关于 $X$ 的超上同调等同。

证明是标准的。考虑双复形 $\Homsh_{\Osh}(L_{*},C(B))$，其中 $C(B)$ 是 $B$
的一个内射消解，再考虑从 $\Homsh_{\Osh}(L_{*},B)$ 和
$\Homsh_{\Osh}(A,C(B))$ 到此双复形的自然同态即可。

最后指出，(1.2) 中引入的两个 $\operatorname{Ext}$ 函子不仅是关于 $B$ 的
上同调函子，而且是上同调双函子：关于 $B$ 协变，关于 $A$ 反变。

\fsection{2}{Ext 中的合成律}

本节结果由 CARTIER 和 YONEDA 各自独立得到；详情见 CARTIER 的报告 [1]。
设 $\Ccat$ 为阿贝尔范畴，$K,L$ 为 $\Ccat$ 中两个分次对象。以
$\operatorname{Hom}(K,L)$ 表示如下分次阿贝尔群：其 $n$ 次分量由从 $K$ 到
$L$ 的 $n$ 次齐次同态组成（即同态 $K^{i}\longrightarrow L^{i+n}$ 所成的
系统 $(u_i)$）。

若 $K,L$ 是复形（为明确起见，设其微分算子的次数为 $+1$），则在
$\operatorname{Hom}(K,L)$ 上引入微分算子
\begin{equation}\tag{2.1}
  \delta(u)=du+(-1)^{n+1}ud,\quad
  \mbox{\upshape\fontspec{SimSun}其中 }n=\deg(u),
\end{equation}
从而使它成为微分次数为 $+1$ 的复形。$n$ 次闭元就是与 $d$ 反交换的 $n$ 次
算子（视为齐次算子）。\footnote{印刷文本作“与 $u$”；方程 (2.1) 要求此处
为“与 $d$”。} 因而可以考虑 $H^{*}\bigl(\operatorname{Hom}(K,L)\bigr)$；
它是 $K,L$ 同伦型的不变量，可记为 $H^{*}(K,L)$。

若还有第三个复形 $M$，则同态的合成定义一个配对
\[
  \operatorname{Hom}(K,L)\times\operatorname{Hom}(L,M)
  \longrightarrow\operatorname{Hom}(K,M)
\]
且与微分算子相容，因此取上同调后得到配对
\begin{equation}\tag{2.2}
  H^{*}(K,L)\times H^{*}(L,M)\longrightarrow H^{*}(K,M),
\end{equation}
记作 $(u,v)\longmapsto vu$。这些配对显然满足结合律。特别地，
$H^{*}(K,K)$ 是带幺结合分次环，而 $H^{*}(K,L)$（相应地，$H^{*}(L,K)$）
是该环上的分次右模（相应地，分次左模），等等。在 $0$ 次时，(2.2) 化为
复形同态的合成。最后，若有复形的正合列
\[
  0\longrightarrow K'\longrightarrow K\longrightarrow K''\longrightarrow 0,
\]

\src{172}
并且每个 $K'^{i}$ 都可与 $K^{i}$ 的一个直和因子等同，则它诱导群复形
$\operatorname{Hom}(K'',L)$ 等的正合列，进而给出连接算子
\[
  H^{i}(K',L)\longrightarrow H^{i+1}(K'',L).
\]
相对于 $L$ 中一条正合列的连接算子也用同样方式定义。配对 (2.2) 按通常意义
与这些连接算子相容。

现在设 $\Ccat$ 是一个范畴，其中每个对象 $A$ 都有内射消解 $C(A)$。运用
双复形定理的诸多变式之一可知
\[
  H^{*}\bigl(C(A),C(B)\bigr)
  =H^{*}\bigl(\operatorname{Hom}(C(A),C(B))\bigr)
\]
典范同构于
\[
  H^{*}\bigl(\operatorname{Hom}(A,C(B))\bigr)
  =\operatorname{Ext}^{*}(A,B).
\]
上面考虑的连接算子给出 $\operatorname{Ext}$ 的连接同态。此外，配对 (2.2)
在此给出与连接同态相容的结合配对：
\begin{equation}\tag{2.3}
  \operatorname{Ext}^{*}(A,B)\times\operatorname{Ext}^{*}(B,C)
  \longrightarrow\operatorname{Ext}^{*}(A,C).
\end{equation}
特别地，$\operatorname{Ext}^{*}(A,A)$ 是带幺结合分次环，等等。（类似地可证，
$\operatorname{Ext}$ 函子作用于任意函子的导出函子；此处不使用这一事实。）

若所考虑的范畴是 $X$ 上 $\Osh$-模的范畴 $\Ccat^{\Osh}$，便得到配对
\begin{equation}\tag{2.4}
  \operatorname{Ext}_{\Osh}^{p}(X;A,B)
  \times\operatorname{Ext}_{\Osh}^{q}(X;B,C)
  \longrightarrow\operatorname{Ext}_{\Osh}^{p+q}(X;A,C),
\end{equation}
它们可以按前述方式明确写出。另一方面，作同样的推导，但把阿贝尔群范畴替换为
$X$ 上阿贝尔群层的范畴，并把 $\operatorname{Hom}$ 函子替换为 $\Homsh$
函子，还可定义具有相同形式性质、而且这一次具有\fquote{局部性质}的配对：
\begin{equation}\tag{2.5}
  \Extsh_{\Osh}^{p}(A,B)\times\Extsh_{\Osh}^{q}(B,C)
  \longrightarrow\Extsh_{\Osh}^{p+q}(A,C).
\end{equation}
再注意同态 (1.8) 与 $\operatorname{Ext}$ 之间的配对相容，便可进一步明确
这些配对。

\src{173}
最后回顾，函子 $H^{p}(X,A)$ 之间也有乘法结构（杯积）。由此可见，命题 1
中的谱序列与乘法结构相容。更确切地说，谱序列 $E(A,B)$ 与谱序列 $E(B,C)$
有一个取值于谱序列 $E(A,C)$ 的配对；在收敛项上，它化为整体
$\operatorname{Ext}$ 群之间的配对，而在初始项上，它化为 (1.4) 右端中由
杯积和局部 $\operatorname{Ext}$ 群配对所导出的配对。特别地，由此可知边缘同态
\begin{equation}\tag{2.6}
  \operatorname{Ext}_{\Osh}^{n}(X;A,B)
  \longrightarrow H^{0}\bigl(X,\Extsh_{\Osh}^{n}(A,B)\bigr)
\end{equation}
和
\begin{equation}\tag{2.7}
  H^{n}\bigl(X,\Homsh_{\Osh}(A,B)\bigr)
  \longrightarrow\operatorname{Ext}_{\Osh}^{n}(X;A,B)
\end{equation}
均与乘法结构相容。因此，若只考虑局部同构于某个 $\Osh^{m}$ 的层，则结合
同构 (1.5)，即可完全用杯积明确描述整体 $\operatorname{Ext}$ 之间的合成。

\fsection{3}{局部上同调的结果}

设 $A$ 是一个带幺交换环，并赋有一个理想 $J$。对变动的 $A$-模 $M$，我们将
定义函子性同态
\begin{equation}\tag{3.1}
\left\{
\begin{aligned}
\operatorname{Ext}_{A}^{p}(A/J,M)
  &\longrightarrow
  \operatorname{Hom}_{A}\!\left(\bigwedge^{p}J/J^{2},M\otimes A/J\right),\\
\operatorname{Tor}_{p}^{A}(A/J,M)
  &\longleftarrow
  \left(\bigwedge^{p}J/J^{2}\right)
  \otimes\operatorname{Hom}_{A}(A/J,M),
\end{aligned}
\right.
\end{equation}
其中张量积和外积都在环 $A$ 上取。还应注意，$J/J^{2}$ 实际上是一个
$A/J$-模，而它作为 $A$-模或作为 $A/J$-模所取的外幂相同。定义同态 (3.1)，
等价于对 $J$ 中元素的每个系统 $\underline{x}=(x_1,\ldots,x_p)$ 定义同态
$\varphi_{\underline{x}}$：
\begin{equation}\tag{3.2}
\left\{
\begin{aligned}
\varphi_{\underline{x}}:
  \operatorname{Ext}_{A}^{p}(A/J,M)&\longrightarrow M\otimes A/J,\\
\varphi_{\underline{x}}:
  \operatorname{Hom}_{A}(A/J,M)&\longrightarrow
  \operatorname{Tor}_{p}^{A}(A/J,M),
\end{aligned}
\right.
\end{equation}
并使其满足下列条件：

\src{174}
\begin{enumerate}
\item[\textup{i.}] $\varphi_{x_1,\ldots,x_p}$ 按交替 $A$-多线性的方式依赖于
$x_i\in J$ 所成的系统。
\item[\textup{ii.}] 只要有一个 $x_i$ 属于 $J^{2}$，就有
$\varphi_{x_1,\ldots,x_p}=0$。
\end{enumerate}
事实上，ii 由 i 推出，因为对 $a\in J$ 有 $a\varphi_{\underline{x}}=0$；只需
注意 (3.2) 中所有的模都被 $J$ 零化。

为定义 $\varphi_{\underline{x}}$，考虑复形 $K_{\underline{x}}$：其底层
$A$-模是 $\bigwedge A^{p}$，其微分算子是内乘算子
$i_{\underline{x}}$，这里将 $\underline{x}$ 视为 $A^{p}$ 上分量为
$x_1,\ldots,x_p$ 的线性形式。微分算子的次数为 $-1$，各次数非负，而该复形
在 $0$ 维的同调为 $A/(x_1A+\cdots+x_pA)$。由于 $x_i\in J$，便得到增广
$K_{\underline{x},0}\longrightarrow A/J$。因此 $K_{\underline{x}}$ 是一个
以 $A/J$ 为增广模的增广自由复形。按熟知的方式，由此得到同态
\[
\operatorname{Ext}_{A}^{*}\bigl(H_0(K_{\underline{x}}),M\bigr)
\longrightarrow
H^{*}\bigl(\operatorname{Hom}_{A}(K_{\underline{x}},M)\bigr)
\]
以及
\[
\operatorname{Tor}_{*}^{A}\bigl(H_0(K_{\underline{x}}),M\bigr)
\longleftarrow
H_{*}\bigl(K_{\underline{x}}\otimes_A M\bigr).
\]
再与增广同态 $H_0(K_{\underline{x}})\longrightarrow A/J$ 所诱导的
$\operatorname{Ext}$ 和 $\operatorname{Tor}$ 上的同态合成，得到同态
\begin{equation}\tag{3.3}
\left\{
\begin{aligned}
\psi_{\underline{x}}:
\operatorname{Ext}_{A}^{*}(A/J,M)
&\longrightarrow H^{*}\bigl(\operatorname{Hom}_{A}(K_{\underline{x}},M)\bigr),\\
\psi_{\underline{x}}:
\operatorname{Tor}_{*}^{A}(A/J,M)
&\longleftarrow H_{*}\bigl(K_{\underline{x}}\otimes_A M\bigr).
\end{aligned}
\right.
\end{equation}
立即可见，在最高维数 $p$，右端的同调是
$M/(x_1M+\cdots+x_pM)$（相应地，是 $M$ 中被每个 $x_i$ 零化的元素所成的
集合）。由于 $x_i\in J$，便得到同态
\begin{equation}\tag{3.4}
\left\{
\begin{aligned}
H^{p}\bigl(\operatorname{Hom}_{A}(K_{\underline{x}},M)\bigr)
&\longrightarrow M\otimes A/J,\\
H_{p}\bigl(K_{\underline{x}}\otimes_A M\bigr)
&\longleftarrow\operatorname{Hom}_{A}(A/J,M).
\end{aligned}
\right.
\end{equation}
将同态 (3.3) 与 (3.4) 合成，便得到所要定义的同态 (3.2)。条件 i 的验证
虽然繁琐，却没有困难。

\src{175}
\result{命题}{4} 设 $A$ 是带幺交换环，$(x_1,\ldots,x_p)$ 是 $A$ 中的一个
元素序列，并且对 $1\leq i\leq p$，$x_i$ 在 $A$ 对由
$(x_1,\ldots,x_{i-1})$ 生成的理想所得商环中的像不是零因子。令 $J$ 为
$x_i$ 生成的理想。则 $J/J^{2}$ 是自由 $A/J$-模，其基为 $x_i$ 的典范像；
复形 $K_{\underline{x}}$ 是 $A/J$ 的一个自由消解；而且对任意 $A$-模 $M$，
(3.1) 中次数为 $p$ 的同态都是双射。只要 $J M=0$，在任意次数 $i$ 定义的
类似同态也具有相同结论。

（关键之处是 $K_{\underline{x}}$ 的无环性；其余结论都由此推出。在上述条件下，
这是熟知的事实。）

\result{推论}{1} 设 $A,J$ 如上，并进一步假设 $A$ 是完美域 $k$ 上维数为
$n$ 的正则仿射代数，且 $A/J$ 也是正则仿射代数。以 $\Omega^{1}(A)$、
$\Omega^{1}(A/J)$ 表示 Kähler 微分模。则有与局部化相容的典范同构
\begin{equation}\tag{3.5}
  \operatorname{Ext}_{A}^{p}\bigl(\Omega^{n-p}(A/J),\Omega^{n}(A)\bigr)=A/J.
\end{equation}
事实上，$\Omega^{n-p}(A/J)$ 是秩为 $1$ 的自由 $A/J$-模，同样，
$\Omega^{n}(A)$ 是秩为 $1$ 的自由 $A$-模，\footnote{印刷文本作“秩 $n$”；
两个最高次形式模的秩均为 $1$，下文与 $\bigwedge^{p}(J/J^{2})$ 的等同也要求
如此。} 因而左端等于
\[
  \operatorname{Ext}_{A}^{p}(A/J,A)
  \otimes\Omega^{n-p}(A/J)'\otimes\Omega^{n}(A)
\]
（其中符号“$'$”表示对偶 $A/J$-模）。后两个因子的张量积可与
$\bigwedge^{p}(J/J^{2})$ 等同，故整个表达式可与
\[
  \operatorname{Ext}_{A}^{p}\!\left(A/J,\bigwedge^{p}(J/J^{2})\right)
\]
等同，再由命题与
\[
  \operatorname{Hom}_{A}\!\left(\bigwedge^{p}J/J^{2},
  \bigwedge^{p}J/J^{2}\right)
\]
等同，也就是与 $A/J$ 等同。特别地，
$\operatorname{Ext}_{A}^{p}\bigl(\Omega^{n-p}(A/J),\Omega^{n}(A)\bigr)$ 中
有一个对应于 $A/J$ 单位元的特出元素，称为\zhterm{理想 $J$ 在 $A$ 中的基本类}。
（事实上，它可以在宽泛得多的条件下定义。）推论 1 可以表为更几何、更整体的
形式：

\result{推论}{2} 设 $X$ 是代数闭域 $k$ 上的非奇异簇，$Y$ 是 $X$ 的闭
非奇异子簇；$\Osh_X$ 是 $X$ 的局部环层，$\Osh_Y$ 是 $Y$ 的局部环层，并将
它视为 $\Osh_X$ 的商层。设 $X$ 的维数为 $n$，$Y$ 的维数为 $n-p$。

\src{176}
令 $\Omega_X$（相应地，$\Omega_Y$）为 $X$（相应地，$Y$）上正则微分形式
的芽层。则有典范同构
\begin{equation}\tag{3.6}
  \Extsh_{\Osh_X}^{p}\bigl(\Omega_Y^{n-p},\Omega_X^{n}\bigr)=\Osh_Y,
\end{equation}
或者等价地，
\begin{equation}\tag{3.6 bis}
  \Extsh_{\Osh_X}^{p}\bigl(\Osh_Y,\Omega_X^{n}\bigr)=\Omega_Y^{n-p}.
\end{equation}
当 $Y$ 是奇异簇时，公式 (3.6 bis) 可用作 $\Omega_Y^{n-p}$ 的定义。更确切地说：

\result{命题}{5} 设 $Y$ 是维数为 $q=n-p$ 的代数子集，包含在维数为 $n$ 的
非奇异代数簇 $X$ 中。设 $F$ 是 $X$ 上支集包含于 $Y$ 的凝聚代数层，$L$ 是
$X$ 上的局部自由代数层。则当 $i<p$ 时，层 $\Extsh_{\Osh_X}^{i}(F,L)$ 为零；
而当 $i=p$ 时，有典范同构
\begin{equation}\tag{3.7}
  \Extsh_{\Osh_X}^{p}(F,L)
  =\Homsh_{\Osh_X}\!\left(F,
  \Extsh_{\Osh_X}^{p}(\Osh_X/J,L)\right),
\end{equation}
其中 $J$ 表示 $X$ 上任意一个零化 $F$、且零点集为 $Y$ 的理想层。特别地，
若 $F$ 是 $Y$ 上的凝聚代数层，则
\begin{equation}\tag{3.7 bis}
  \Extsh_{\Osh_X}^{p}(F,L)
  =\Homsh_{\Osh_Y}\!\left(F,\Extsh_{\Osh_X}^{p}(\Osh_Y,L)\right).
\end{equation}
最后，仍设 $F$ 是 $Y$ 上的凝聚代数层，则诸层
\[
  E^{i}(F)=\Extsh_{\Osh_X}^{p+i}(F,\Omega_X^{n})
\]
不依赖于把代数空间 $Y$ 嵌入非奇异簇 $X$ 的方式。

问题是局部的，故可设 $X$ 仿射且 $L=\Osh_X$。这把问题归结为交换代数问题，
甚至局部代数问题：若 $A$ 是正则局部环，$M$ 是支集维数
$\leq q=n-p$ 的 $A$-模，则要证明当 $i<p$ 时
$\operatorname{Ext}_{A}^{i}(M,A)=0$，并且
\[
  \operatorname{Ext}_{A}^{p}(M,A)
  =\operatorname{Hom}_{A}\!\left(M,\operatorname{Ext}_{A}^{p}(A/J,A)\right),
\]
其中 $J$ 是任意一个零化 $M$ 且\fquote{维数} $\leq q$ 的理想。对于第一点，
对 $q$ 作归纳：一次直接的分解化简把问题归约到 $M=A/J$ 的情形；随后把 $J$
换成一个更小的理想，并运用归纳假设和 $\operatorname{Ext}$ 的正合列，又可归约
到 $J$ 由命题 4 中那样的一个\fquote{参数系}生成的情形，此时结论立即成立。
前一结果蕴含：若 $J$ 是一个理想

\src{177}
固定且\fquote{维数} $\leq q$，则 $A/J$-模范畴上的反变函子
$E(M)=\operatorname{Ext}_{A}^{p}(M,A)$ 是左正合的；它还把直和变为直积，
由此容易得到
\[
  E(M)=\operatorname{Hom}_{A}\bigl(M,E(A/J)\bigr).
\]
\footnote{印刷文本作 $E(A)$；在所考虑的范畴中，表示对象是 $A/J$，而 (3.7)
所需的等式为 $E(A/J)$。} 最后，命题 5 的最后一个断言更为微妙；它由
$E^{i}(F)$ 的一个内蕴刻画推出，而该刻画使用了此处无法陈述的局部对偶定理。

\result{推论}{} 以 $\omega_Y^{q}$ 表示层
$\Extsh_{\Osh_X}^{p}(\Osh_Y,\Omega_X^{n})$。则对 $Y$ 上的凝聚代数层 $F$，
有函子性同构
\begin{equation}\tag{3.8}
  \Extsh_{\Osh_X}^{p}(F,\Omega_X^{n})
  =\Homsh_{\Osh_X}(F,\omega_Y^{q}).
\end{equation}

\fsection{4}{与子簇相伴的上同调类}

下文始终以 $X$ 表示定义在域 $k$ 上的维数为 $n$ 的代数集；为简便起见，假设
$k$ 代数闭。除第 6 节外，均假设 $X$ 非奇异。以 $\Osh_X$ 表示 $X$ 上的局部
环层，以
\[
  \Omega_X^{*}=\bigoplus_p\Omega_X^{p}
\]
表示 $X$ 上微分形式的芽层。\footnote{印刷文本使用并集号；这里的分次结构是
诸层 $\Omega_X^p$ 的直和。} 若 $Y$ 是 $X$ 的闭子集，就把 $Y$ 上的凝聚代数层
视同为 $X$ 上在 $Y$ 外为零的凝聚层；$\Osh_Y$ 和 $\Omega_Y$ 尤其如此。

\result{引理}{1} 设 $F$ 是 $X$ 上支集维数 $\leq n-p$ 的凝聚代数层，$L$
是 $X$ 上的局部自由凝聚代数层。则当 $i<p$ 时，
$\operatorname{Ext}_{\Osh_X}^{i}(X;F,L)$ 为零，并有典范同构
\begin{equation}\tag{4.1}
  \operatorname{Ext}_{\Osh_X}^{p}(X;F,L)
  =H^{0}\bigl(X,\Extsh_{\Osh_X}^{p}(F,L)\bigr).
\end{equation}
若 $F$ 是 $X$ 的闭子集 $Y$ 上的凝聚代数层，且 $Y$ 的维数
$\leq n-p$，则有典范同构
\begin{equation}\tag{4.1 bis}
  \Extsh_{\Osh_X}^{p}(F,L)
  =\Homsh_{\Osh_X}\!\left(
    F\otimes L'\otimes\Omega_X^{n},\omega_Y^{n-p}
  \right),
\end{equation}
其中 $\omega_Y^{n-p}$ 是命题 5 之后的推论所定义的 $Y$ 上的层（若 $Y$ 非奇异，
它可与 $\Omega_Y^{n-p}$ 等同）。

公式 (4.1) 直接由命题 1 的谱序列和命题 5 推出；根据公式 (3.8)，可以写成

\src{178}
\begin{align*}
\Extsh_{\Osh_X}^{p}(F,L)
&=L\otimes(\Omega_X^{n})'\otimes
  \Extsh_{\Osh_X}^{p}(F,\Omega_X^{n})\\
&=L\otimes(\Omega_X^{n})'\otimes
  \Homsh_{\Osh_X}(F,\omega_Y^{q})\\
&=\Homsh_{\Osh_X}\!\left(
  F\otimes L'\otimes\Omega_X^{n},\omega_Y^{q}\right),
\end{align*}
其中 $q=n-p$，由此得到公式 (4.1 bis)。

特别取 $F=\Osh_Y$、$L=\Omega_X^{p}$，并注意到
$\Omega_X^{n}\otimes(\Omega_X^{p})'=\Omega_X^{n-p}$，便得到典范同构
\begin{equation}\tag{4.2}
  \operatorname{Ext}_{\Osh_X}^{p}(X;\Osh_Y,\Omega_X^{p})
  =\operatorname{Hom}_{\Osh_X}
  (X;\Omega_X^{n-p},\omega_Y^{n-p}).
\end{equation}
为简便起见，现在假设 $Y$ 非奇异，因而
$\omega_Y^{n-p}=\Omega_Y^{n-p}$。从 $\Omega_X^{n-p}$ 到
$\Omega_Y^{n-p}$ 有自然同态，故得到层
$\Extsh_{\Osh_X}^{p}(\Osh_Y,\Omega_X^{p})$ 的典范截面 $s_Y$。若 $Y$ 的
所有分支维数均为 $n-p$，就称它为层
$\Extsh_{\Osh_X}^{p}(\Osh_Y,\Omega_X^{p})$ 的\zhterm{基本截面}。由 (4.1)，
该截面定义 $\operatorname{Ext}_{\Osh_X}^{p}(X;\Osh_Y,\Omega_X^{p})$ 中的
一个元素。另一方面，自然同态 $\Osh_X\longrightarrow\Osh_Y$ 定义同态
\[
  \operatorname{Ext}_{\Osh_X}^{p}(X;\Osh_Y,\Omega_X^{p})
  \longrightarrow
  \operatorname{Ext}_{\Osh_X}^{p}(X;\Osh_X,\Omega_X^{p})
  =H^{p}(X,\Omega_X^{p}).
\]
这样便得到 $H^{p}(X,\Omega_X^{p})$ 中的一个元素，记作 $P_X(Y)$，称为
\zhterm{$Y$ 在 $X$ 中的上同调类}。因此，它由
$\Extsh_{\Osh_X}^{p}(\Osh_Y,\Omega_X^{p})$ 的截面 $s_Y$ 通过下列同态图式
得到：
\begin{equation}\tag{4.3}
\begin{aligned}
H^{p}(X,\Omega_X^{p})
&=\operatorname{Ext}_{\Osh_X}^{p}(X;\Osh_X,\Omega_X^{p})\\
&\longleftarrow\operatorname{Ext}_{\Osh_X}^{p}(X;\Osh_Y,\Omega_X^{p})
\xrightarrow{\ \sim\ }
H^{0}\!\left(X,\Extsh_{\Osh_X}^{p}(\Osh_Y,\Omega_X^{p})\right).
\end{aligned}
\end{equation}

所谓\zhterm{维数为 $n-p$ 的非奇异循环}，是指由 $X$ 中维数为 $n-p$ 的不可约
非奇异子簇生成的自由阿贝尔群中的元素。于是函数 $Y\longmapsto P(Y)$ 延拓为
从 $X$ 上维数为 $n-p$ 的非奇异循环群到群 $H^{p}(X,\Omega_X^{p})$ 的同态。

设 $Z^{n-p}$ 和 $Z'^{n-p'}$ 是维数分别为 $n-p$ 和 $n-p'$ 的非奇异循环。
若 $Z$ 的每个分支都与 $Z'$ 的每个分支横截相交，就称它们\zhterm{横截相交}。
此时循环 $Z.Z'$ 有定义，并且是维数为 $n-p-p'$ 的非奇异循环。在这些约定下，
有如下结论。

\src{179}
\result{定理}{1} 若 $Z^{n-p}$ 和 $Z'^{n-p'}$ 是横截相交的非奇异循环，则
\begin{equation}\tag{4.4}
  P_X(Z.Z')=P_X(Z).P_X(Z'),
\end{equation}
其中右端的乘法是杯积
\[
  H^{p}(X,\Omega_X^{p})\times H^{p'}(X,\Omega_X^{p'})
  \longrightarrow H^{p+p'}(X,\Omega_X^{p+p'}).
\]
（假设 $X$ 同构于某个射影空间的一个局部闭子集。）

最后这个假设仅用于保证 $X$ 上每个凝聚代数层都是某个局部自由凝聚代数层的
商（SERRE），因而具有由局部自由层组成的左消解。为证明定理 1，可以假设
$Z,Z'$ 是横截相交的不可约非奇异子簇 $Y,Y'$。设 $L_{*}$ 是 $\Osh_Y$ 的一个
由局部自由层组成的左消解；则根据命题 3，同态图式 (4.3) 可与下图等同：
\[
R^{p}\Gamma(\Omega_X^{p})
\xleftarrow{\alpha}
\mathbb R^{p}\Gamma\!\left(\Homsh_{\Osh_X}(L_{*},\Omega_X^{p})\right)
\xrightarrow[\sim]{\beta}
\Gamma\!\left(H^{p}\bigl(\Homsh_{\Osh_X}(L_{*},\Omega_X^{p})\bigr)\right).
\]
这里 $\Gamma$ 是 $X$ 上阿贝尔群层范畴的\fquote{截面群}函子，
$\mathbb R^{p}\Gamma$ 表示该函子的 $p$ 维超上同调，而 $R^{p}\Gamma$ 表示其
第 $p$ 个导出函子。为简便起见，可设 $L_0=\Osh_X$，并设增广
$L_0\longrightarrow\Osh_Y$ 为自然同态（这是允许的）。此时 $\alpha$ 由复形
同态 $\Osh_X\longrightarrow L_{*}$ 诱导（把 $\Osh_X$ 视为集中在次数 $0$ 的
复形）；这里用到：若 $K$ 是集中在次数 $0$ 的层复形，则
$\mathbb R^{p}\Gamma(K)=R^{p}\Gamma(K_0)$。同态 $\beta$ 是熟知的边缘同态。
对 $\Osh_{Y'}$ 的一个局部自由消解 $L'_{*}$ 考虑类似图式，再考虑下列配对的
交换图：
\begin{equation}\tag{4.5}
\resizebox{\textwidth}{!}{$
\begin{array}{ccc}
R^{p}\Gamma(\Omega_X^{p})
&\longleftarrow
\mathbb R^{p}\Gamma\!\left(\Homsh_{\Osh_X}(L_{*},\Omega_X^{p})\right)
&\xrightarrow{\sim}
\Gamma H^{p}\!\left(\Homsh_{\Osh_X}(L_{*},\Omega_X^{p})\right)
\\[3pt]
\times&\times&\times\\[3pt]
R^{p'}\Gamma(\Omega_X^{p'})
&\longleftarrow
\mathbb R^{p'}\Gamma\!\left(\Homsh_{\Osh_X}(L'_{*},\Omega_X^{p'})\right)
&\xrightarrow{\sim}
\Gamma H^{p'}\!\left(\Homsh_{\Osh_X}(L'_{*},\Omega_X^{p'})\right)
\\[3pt]
\downarrow&\downarrow&\downarrow\\[3pt]
R^{p+p'}\Gamma(\Omega_X^{p+p'})
&\longleftarrow
\mathbb R^{p+p'}\Gamma\!\left(
\Homsh_{\Osh_X}(L_{*}\otimes L'_{*},\Omega_X^{p+p'})\right)
&\xrightarrow{\sim}
\Gamma H^{p+p'}\!\left(
\Homsh_{\Osh_X}(L_{*}\otimes L'_{*},\Omega_X^{p+p'})\right).
\end{array}
$}
\end{equation}
\src{180}
(4.5) 右侧两列中的配对由下列层复形的配对诱导：
\[
\Homsh_{\Osh_X}(L_{*},\Omega_X^{p})
\times\Homsh_{\Osh_X}(L'_{*},\Omega_X^{p'})
\longrightarrow
\Homsh_{\Osh_X}(L_{*}\otimes L'_{*},\Omega_X^{p+p'}),
\]
该配对借助外积
$\Omega_X^{p}\times\Omega_X^{p'}\longrightarrow\Omega_X^{p+p'}$ 定义。
第一列中的配对是（相对于此外积的）杯积。我们断言，(4.5) 的最后一行可与
类似于 (4.3) 的同构图式等同，其中以 $Y\cap Y'$ 代替 $Y$，以 $p+p'$ 代替
$p$。为此只需证明 $L\otimes L'$ 是 $\Osh_{Y\cap Y'}$ 的一个消解（它显然
局部自由）。事实上，
\[
  H_0(L\otimes L')=\Osh_Y\otimes\Osh_{Y'}=\Osh_{Y\cap Y'}
\]
且
\[
  H_i(L\otimes L')=\operatorname{Tor}_{i}^{\Osh_X}(\Osh_Y,\Osh_{Y'})=0
  \qquad(i>0),
\]
因为 $Y,Y'$ 横截相交。现在定理 1 由下式推出：
\begin{equation}\tag{4.6}
  s_Y.s_{Y'}=s_{Y.Y'}
\end{equation}
（左端的乘法是 (4.5) 最后一列中的乘法）。公式 (4.6) 具有纯局部性质；对
$L_{*},L'_{*}$ 取命题 4 中所考虑的消解，即可毫无困难地验证它。类似地（而且
证明更容易）可证 $Z\longmapsto P_X(Z)$ 与笛卡儿积相容：
\begin{equation}\tag{4.7}
  P_{X\times X'}(Z\times Z')=P_X(Z)\otimes P_{X'}(Z')
\end{equation}
（当 $Z$、相应地 $Z'$ 是非奇异簇 $X$、相应地 $X'$ 上的非奇异循环时，该式
成立；这里把 $Z\times Z'$ 视为 $X\times X'$ 上的非奇异循环）。由 (4.4)
和 (4.7) 可知，$P_X(Z)$ 还与非奇异簇之间的态射
$f:X\longrightarrow X'$ 所给出的\fquote{逆像}运算相容：
\begin{equation}\tag{4.8}
  P_X\bigl(f^{-1}(Z')\bigr)=f^{*}\bigl(P_{X'}(Z')\bigr),
\end{equation}
当 $Z'$ 是 $X'$ 上的非奇异循环且 $f$ 与 $Z'$\fquote{横截}时，该式成立；
也就是说，$f$ 的图像在 $X\times X'$ 中与循环 $X\times Z'$ 横截。

\src{181}
\result{推论}{1} 设 $X,X'$ 是某个射影空间中的两个局部闭非奇异簇，并设
$X'$ 完备。设 $U$ 是 $X\times X'$ 上的非奇异循环，$a,b$ 是 $X'$ 中的两个
点，并且 $U$ 与循环 $X\times(a)$ 和 $X\times(b)$ 都横截相交。设 $Z,Z'$ 是
$X$ 上满足
\[
  Z\times(a)=(X\times(a)).U,
  \qquad
  Z'\times(b)=(X\times(b)).U
\]
的非奇异循环。则
\[
  P_X(Z)=P_X(Z').
\]

事实上，令 $f_a:X\longrightarrow X\times X'$ 由 $f_a(x)=(x,a)$ 定义，
根据 (4.8) 有 $P(Z)=f_a^{*}(P(U))$；同理，$P(Z')=f_b^{*}(P(U))$。利用
Künneth 公式
\[
  H^{*}(X\times X',\Omega_{X\times X'}^{*})
  =H^{*}(X,\Omega_X^{*})\otimes H^{*}(X',\Omega_{X'}^{*})
\]
以及 $H^{0}(X',\Osh_{X'})$ 仅由标量组成这一事实，容易得到
$f_a^{*}=f_b^{*}$，从而得证。

对每个 $x\in X$，点 $(x)$ 都是 $X$ 中余维为 $n$ 的非奇异子簇，因而定义
$H^{n}(X,\Omega_X^{n})$ 中的一个元素 $\varepsilon_x$。若 $X$ 是非奇异射影簇，
由推论 1 可知 $\varepsilon_x$ 不依赖于所选的点 $x$；将其记为
$\varepsilon_X$，称为 $H^{n}(X,\Omega_X^{n})$ 的\zhterm{基本类}。

\result{注}{} 为得到令人满意的理论，应当对任意循环 $Z$ 定义 $P_X(Z)$，并对
循环的正常相交证明定理 1。（撰写本报告时，这一点尚未在完全一般的情形中完成。）
假定这些工作已经完成，推论 1 就变成：若 $Z,Z'$ 是两个代数等价的循环，则
$P_X(Z)=P_X(Z')$（即使 $Z,Z'$ 非奇异，这一陈述似乎也不能由前述结果推出）。

\fsection{5}{对偶定理}

本节中，$X$ 表示维数为 $n$ 的非奇异射影簇。

\result{定理}{2} $H^{n}(X,\Omega_X^{n})$ 的基本类 $\varepsilon_X$ 是该向量
空间的一组基。

（证明将在下文给出。）借助前一定理，可以把 $H^{n}(X,\Omega_X^{n})$ 与域
$k$ 等同。现在考虑第 2 节中引入的配对；它特别给出配对

\src{182}
\[
\operatorname{Ext}_{\Osh_X}^{p}(X;\Osh_X,F)
\times\operatorname{Ext}_{\Osh_X}^{n-p}(X;F,\Omega_X^{n})
\longrightarrow
\operatorname{Ext}_{\Osh_X}^{n}(X;\Osh_X,\Omega_X^{n}),
\]
即
\begin{equation}\tag{5.1}
  H^{p}(X,F)
  \times\operatorname{Ext}_{\Osh_X}^{n-p}(X;F,\Omega_X^{n})
  \longrightarrow H^{n}(X,\Omega_X^{n}).
\end{equation}
考虑到定理 2，该配对定义同态
\begin{equation}\tag{5.2}
  \operatorname{Ext}_{\Osh_X}^{n-p}(X;F,\Omega_X^{n})
  \longrightarrow\bigl(H^{p}(X,F)\bigr)'.
\end{equation}
该同态关于 $F$ 具有函子性，并与正合列
$0\longrightarrow F'\longrightarrow F\longrightarrow F''\longrightarrow0$
所对应的连接同态可交换。

\result{定理}{3} 同态 (5.2) 是同构。

特别地，由此重新得到 Serre 的结果：

\result{推论}{} 设 $E$ 是 $X$ 上的代数向量丛，$\Osh_X(E)$ 是 $E$ 的正则
截面芽层。则有典范同构
\begin{equation}\tag{5.3}
  \bigl(H^{p}(X,\Osh_X(E))\bigr)'
  =H^{n-p}\bigl(X,\Omega_X^{n}\otimes\Osh_X(E')\bigr).
\end{equation}
只需把定理 3 和命题 1 的推论 1 结合使用。

定理 2 和定理 3 将由下述命题推出：

\noindent\textbf{(D) 下列同态}
\begin{equation}\tag{5.2 bis}
  \operatorname{Ext}_{\Osh_X}^{n-p}(X;F,\Omega_X^{n})
  \longrightarrow\bigl(H^{p}(X,F)\bigr)'\otimes L,
  \qquad L=H^{n}(X,\Omega_X^{n}),
\end{equation}
由配对 (5.1) 诱导，并且是同构。

先证明 (D) 蕴含定理 2。令 $k_x=\Osh_{(x)}$ 为缩约到点 $x\in X$ 的簇的
局部环层。考虑典范同态 $\Osh_X\longrightarrow k_x$ 以及相伴的同态
\begin{equation}\tag{5.3 bis}
  H^{0}(X,\Osh_X)\longrightarrow H^{0}(X,k_x),
\end{equation}
其转置可与下列同态等同：
\begin{equation}\tag{5.4}
  \operatorname{Ext}_{\Osh_X}^{n}(X;k_x,\Omega_X^{n})\otimes L'
  \longrightarrow
  \operatorname{Ext}_{\Osh_X}^{n}(X;\Osh_X,\Omega_X^{n})\otimes L'
\end{equation}
该同态由 $\Osh_X\longrightarrow k_x$ 所对应的 $\operatorname{Ext}^{n}$ 上的
同态
\begin{equation}\tag{5.5}
  \operatorname{Ext}_{\Osh_X}^{n}(X;k_x,\Omega_X^{n})
  \longrightarrow
  \operatorname{Ext}_{\Osh_X}^{n}(X;\Osh_X,\Omega_X^{n})
  =H^{n}(X,\Omega_X^{n}).
\end{equation}
诱导。

\src{183}
由于 (5.3 bis) 是同构，\footnote{印刷文本在第 182 页给出了第二个公式
“(5.3)”，并在此处引用它；标识“(5.3 bis)”把该公式与第一个 (5.3) 区分开。}
(5.4) 也是同构，因而 (5.5) 也是同构。根据 (4.2)，$s_{(x)}$ 是
$\operatorname{Ext}_{\Osh_X}^{n}(X;k_x,\Omega_X^{n})$ 的一组基，故其像
$\varepsilon_X$ 确实是 $H^{n}(X,\Omega_X^{n})$ 的一组基。

还需证明命题 (D)。它将由以下诸引理所概括的初等事实纯形式地推出。在这些引理
中，假设 $X$ 是维数为 $r$ 的射影空间 $P$ 的闭子集（可以奇异，也可以非奇异）。
以 $\Osh_P(m)$ 表示 [3] 中记作 $\Osh(m)$ 的 $P$ 上的层，以 $\Osh_X(m)$
表示 $X$ 上类似的层。

\result{引理}{2} 若 $X=P$ 且 $F=\Osh_P(m)$，则命题 (D) 成立。

该引理可由直接计算验证。$H^{i}(P,\Osh_P(m))$ 的显式计算见 [3]，不过也可用
更初等的方法完成。为明确写出配对 (5.1) 而需要计算的杯积
\[
  H^{i}(P,\Osh_P(m))\times H^{j}(P,\Osh_P(m'))
  \longrightarrow H^{i+j}(P,\Osh_P(m+m'))
\]
没有困难。\footnote{印刷文本作 $\Osh_P(r+r')$；两个因子 $\Osh_P(m)$ 和
$\Osh_P(m')$ 要求此处为 $\Osh_P(m+m')$。}

\result{引理}{3} $X$ 上每个凝聚代数层 $F$ 都同构于某个层
$\Osh_X(-m)^{k}$ 的商层，其中 $m$ 可以取得任意大。

这是因为当 $m$ 足够大时，$F\otimes\Osh_X(m)$\fquote{由其截面生成}；
参见 [3]。

\result{引理}{4} 设 $i>0$。当 $m$ 足够大时，
$H^{r-i}(P,\Osh_P(-m))=0$；并且对 $X$ 上任意凝聚代数层 $B$，当 $m$
足够大时有
\[
  \operatorname{Ext}_{\Osh_X}^{i}(X;\Osh_X(-m),B)=0.
\]

第一个断言由上面提到的显式计算推出。对第二个断言，注意有同构
\[
  \operatorname{Ext}_{\Osh_X}^{i}(X;\Osh_X(-m),B)
  =H^{i}(X,B\otimes\Osh_X(m))
\]
（命题 1，推论 1），再由 [3] 中一个熟知的结果即得结论。结合前两个引理，
得到以下推论。

\result{推论}{} 设 $i>0$。则 $P$ 上凝聚代数层范畴中的函子
$F\longmapsto H^{r-i}(P,F)$ 是余可抹去的；$X$ 上凝聚代数层范畴中的函子
$\operatorname{Ext}_{\Osh_X}^{i}(X;F,B)$ 也是余可抹去的。

\src{184}
\result{引理}{5} 设 $A,B$ 是 $X$ 上两个凝聚代数层，并令
$A(m)=A\otimes\Osh_X(m)$。则当 $m$ 足够大时，典范同态
\[
  \operatorname{Ext}_{\Osh_X}^{1}(X;A(-m),B)
  \longrightarrow H^{0}\bigl(X,\Extsh_{\Osh_X}^{1}(A(-m),B)\bigr)
  =H^{0}\bigl(X,\Extsh_{\Osh_X}^{1}(A,B)(m)\bigr)
\]
是同构。

把命题 1 的谱序列用于 $A(-m),B$，便立即得到此结论，因为此时
\[
  E_{2}^{p,q}(A(-m),B)
  =H^{p}\bigl(X,\Extsh_{\Osh_X}^{q}(A(-m),B)\bigr)
  =H^{p}\bigl(X,\Extsh_{\Osh_X}^{q}(A,B)(m)\bigr),
\]
而当 $p>0$ 且 $m$ 足够大时，此式为零。

现在证明 $X=P$ 时的 (D)。先证当 $p=n$ 时 (5.2 bis) 为同构；此时两端都是
左正合函子（因为 $H^{r+1}(P,F)=0$），故由引理 3，只需对
$F=\Osh_P(-m)$ 证明该断言，而这正包含在引理 2 中。由于同态 (5.2 bis)
具有函子性并与连接算子相容，且当 $p<n$ 时 (5.2 bis) 两端都是关于 $F$ 的
余可抹去函子（引理 4 的推论），由标准论证可知 (5.2 bis) 对每个 $p$ 都是
同构。这就证明了射影空间的对偶定理。

现在设 $X$ 任意但非奇异。由 $P$ 的对偶定理，有同构
\[
  H^{n}(X,F)'=H^{n}(P,F)'
  =\operatorname{Ext}_{\Osh_P}^{r-n}(P;F,\Omega_P^{r}).
\]
根据第 4 节引理 1，最后一端可与
\[
  \operatorname{Hom}_{\Osh_P}(P;F,\omega_X^{n})
  =\operatorname{Hom}_{\Osh_X}(X;F,\Omega_X^{n})
  =\operatorname{Ext}_{\Osh_X}^{0}(X;F,\Omega_X^{n})
\]
等同，因而有同构
\begin{equation}\tag{5.6}
  H^{n}(X,F)'=\operatorname{Hom}_{\Osh_X}(X;F,\Omega_X^{n})
  =\operatorname{Ext}_{\Osh_X}^{0}(X;F,\Omega_X^{n}).
\end{equation}
取 $F=\Omega_X^{n}$，得到同构
\begin{equation}\tag{5.7}
  \gamma:H^{n}(X,\Omega_X^{n})\xrightarrow{\sim}k.
\end{equation}

\src{185}
可以验证，借助 (5.7) 取 $L=k$ 后，同构 (5.6) 正是 $p=n$ 时的
(5.2 bis)。因此当 $p=n$ 时，(5.2 bis) 是同构。为证明它对每个 $p$ 都是
同构，只需再证明：当 $p<n$ 时，(5.2 bis) 两端都是关于 $F$ 的余可抹去函子；
进一步说（考虑到引理 3），只需证明取 $F=\Osh_X(-m)$ 且 $m$ 足够大时两端
都为零。对第一端，这由引理 4 成立；对第二端，利用 $P$ 上的对偶定理写成
\[
  H^{p}(X,\Osh_X(-m))'
  =\operatorname{Ext}_{\Osh_P}^{r-p}
    (P;\Osh_X(-m),\Omega_P^{r}).
\]
当 $p<n$ 且 $m$ 足够大时，第二端为零。这由引理 5（其中取 $X=P$）以及下述
事实推出：$\Osh_X$ 作为 $P$ 上的凝聚代数层，其上同调维数
$\leq r-n$（因为 $X$ 非奇异），所以
\[
  \Extsh_{\Osh_P}^{r-p}(\Osh_X,\Omega_P^{r})=0
  \qquad\mbox{\upshape\fontspec{SimSun}当 }p<n.
\]

\fsection{6}{奇异簇的对偶定理}

设 $X$ 是维数为 $r$ 的射影空间 $P$ 中维数为 $n$ 的闭子集。此处公式 (5.6)
应写成
\begin{equation}\tag{6.1}
  H^{n}(X,F)'
  \simeq\operatorname{Hom}_{\Osh_X}(X;F,\omega_X^{n})
  =\operatorname{Ext}_{\Osh_X}^{0}(X;F,\omega_X^{n}),
\end{equation}
其中定义
\begin{equation}\tag{6.2}
  \omega_X^{n}=E^{0}(\Osh_X)
  =\Extsh_{\Osh_P}^{r-n}(\Osh_X,\Omega_P^{r}).
\end{equation}
如命题 5 所述，这样引入的层实际上不依赖于所选的把 $X$ 嵌入非奇异簇 $P$
的方式。在 (6.1) 中取 $F=\omega_X^{n}$，得到
\begin{equation}\tag{6.2 bis}
  H^{n}(X,\omega_X^{n})'
  \simeq\operatorname{Hom}_{\Osh_X}(X;\omega_X^{n},\omega_X^{n}),
\end{equation}
\footnote{印刷文本也把该公式编号为“(6.2)”；后缀“bis”用于区分两次出现。}
因而 $H^{n}(X,\omega_X^{n})'$ 中有一个对应于 $\omega_X^{n}$ 恒等同态的
特出元素：
\begin{equation}\tag{6.3}
  \gamma:H^{n}(X,\omega_X^{n})\longrightarrow k.
\end{equation}

\src{186}
现考察由 $\operatorname{Ext}$ 的复合所定义的配对：
\begin{equation}\tag{6.4}
  H^{p}(X,F)\times
  \operatorname{Ext}_{\Osh_X}^{n-p}(X;F,\omega_X^{n})
  \longrightarrow H^{n}(X,\omega_X^{n}),
\end{equation}
再把它们与 (6.3) 中的同态 $\gamma$ 复合，便得到与连接算子相容的函子性同态
\begin{equation}\tag{6.5}
  \operatorname{Ext}_{\Osh_X}^{n-p}(X;F,\omega_X^{n})
  \longrightarrow H^{p}(X,F)'
\end{equation}
（它推广了 (5.2)）。可以验证，当 $p=n$ 时，由此得到同构 (6.1)。在这些约定下，
有如下结论。

\result{定理}{3 bis} 对给定整数 $k\geq0$，下列关于 $X$ 的四个条件等价：
\begin{enumerate}
  \item[\textit{i.}] 当 $n-k\leq p\leq n$ 时，函子性同态 (6.5) 是同构。
  \item[\textit{ii.}] 当 $m$ 足够大且 $n-k\leq p<n$ 时，
    $H^{p}(X,\Osh_X(-m))=0$。
  \item[\textit{iii.}] 当 $n-k\leq p<n$ 时，$X$ 上凝聚代数层范畴中的函子
    $H^{p}(X,F)$ 是余可抹去的。
  \item[\textit{iv.}] 当 $0<i\leq k$ 时，
    $E^{i}(\Osh_X)=\Extsh_{\Osh_P}^{r-n+i}(\Osh_X,\Omega_P^{r})=0$。
\end{enumerate}

\noindent\textbf{证明。}——由引理 4 得
\textit{i.}$\Rightarrow$\textit{ii.}，由引理 3 得
\textit{ii.}$\Rightarrow$\textit{iii.}。又因为当 $n-k\leq p<n$ 时，
(6.5) 两端都是余可抹去函子（前者由引理 4），故由一个熟知的标准论证得到
\textit{iii.}$\Rightarrow$\textit{i.}。最后证明
\textit{ii.}$\Leftrightarrow$\textit{iv.}；这由下一命题的推论推出。

\result{命题}{6} 设 $F$ 是 $X$ 上的凝聚代数层，$i$ 是整数。则当 $m$ 足够大时，
有同构
\begin{equation}\tag{6.6}
  H^{i}(X,F(-m))'
  \simeq H^{0}\bigl(X,E^{n-i}(F)(m)\bigr),
\end{equation}
其中规定（比较第 3 节命题 5）
\begin{equation}\tag{6.7}
  E^{j}(F)=\Extsh_{\Osh_P}^{r-n+j}(F,\Omega_P^{r}).
\end{equation}
事实上，由 $P$ 上的对偶定理，(6.6) 的左端同构于
$\operatorname{Ext}_{\Osh_P}^{r-i}(P;F(-m),\Omega_P^{r})$，故 (6.6)
由引理 5 推出。

\src{187}
\result{推论}{} 当 $m$ 足够大时 $H^{i}(X,F(-m))=0$，当且仅当
$E^{n-i}(F)=0$。

回顾 $E^{j}(F)$ 实际上不依赖于所考察的射影嵌入。该推论中的条件完全是局部的；
因此，若它对 $F$ 成立，就对任一局部同构于层 $F^{n}$ 的层成立。特别地，若该
条件对 $\Osh_X$ 成立，就对每个局部自由凝聚代数层成立。例如：若 $X$ 非奇异，
则对每个 $i<n$ 都成立；若 $X$ 的任何分支都不退化为一个点，则对 $i=0$
成立；若 $X$ 正规且其所有分支的维数都 $>1$，则对 $i=0,1$ 成立（参见
[3]）。\footnote{印刷文本此处写作“$S$ 正规”；本段只涉及代数簇 $X$。}
按照定义，要使它对每个 $i<k$ 都成立，充要条件是局部环
$\Osh_x$（$x\in X$）的\fquote{同调余维数} $\geq k$（关于此概念参见
[4]）。若 $k=n$，则由定理 3 bis，这意味着对偶定理对 $X$ 成立，即对每个
$p$ 和每个 $F$，(6.5) 都是同构。关于局部环 $\Osh_x$，可以给出许多与此
等价的条件（NAGATA），例如要求它们满足 Cohen--Macaulay 等维定理。比如，
若 $X$ 局部地是某个非奇异环境代数簇中的\fquote{完全交}，这一条件便成立。

\fsection{7}{庞加莱对偶}

设 $X$ 是维数为 $n$ 的非奇异射影代数簇。则
\[
  H^{**}(X)=H^{*}(X,\Omega_X^{*})
\]
是一个有限维反交换双分次代数；我们按总次数给它分次，因而
$H^{p,q}(X)=H^{p}(X,\Omega_X^{q})$ 的次数为 $p+q$。$H^{*}(X)$ 的次数介于
$0$ 与 $2n$ 之间。由定理 2 和定理 3 的推论，$H^{*}(X)$ 是维数为 $2n$ 的
\fquote{庞加莱代数}：也就是说，$H^{2n}(X)$ 具有到基域 $k$ 的一个同构，且乘法
\[
  H^{m}(X)\times H^{2n-m}(X)\longrightarrow H^{2n}(X)=k
\]
给出 $H^{m}(X)$ 与 $H^{2n-m}(X)$ 之间的对偶。此外，若 $Y$ 是另一个非奇异
射影代数簇，则凝聚代数层的 Künneth 公式给出
\begin{equation}\tag{7.1}
  H^{*}(X\times Y)=H^{*}(X)\otimes H^{*}(Y),
\end{equation}
该同构与庞加莱代数结构相容。再者，作为交换代数，$H^{*}(X)$ 是 $X$ 的反变
函子；态射 $f:Y\longrightarrow X$ 以显然方式定义分次代数同态

\src{188}
\begin{equation}\tag{7.2}
  f^{*}:H^{*}(X)\longrightarrow H^{*}(Y).
\end{equation}
由于这里涉及的是庞加莱代数，转置便给出向量空间同态
\begin{equation}\tag{7.3}
  f_{*}:H^{*}(Y)\longrightarrow H^{*}(X).
\end{equation}
我们在第 4 节已经看到，$f^{*}$ 对非奇异循环所对应上同调类的作用，可在几何上
解释为取这些循环的逆像所对应的上同调类。在当前情形，重要的是证明 (7.3)
同样对应于循环的\fquote{正像}运算。至少在适当的非奇异性条件下，这由下述特殊
情形推出。

\result{定理}{4} 若 $f$ 是非奇异子簇 $Y^{m}\subset X^{n}$ 到 $X^{n}$ 的
自然包含映射，并以 $1_Y$ 表示 $H(Y)$ 的单位元，则
\begin{equation}\tag{7.4}
  f_{*}(1_Y)=P_X(Y),
\end{equation}
其中右端是 $Y$ 在 $X$ 中所对应的上同调类。

该公式等价于
\begin{equation}\tag{7.4 bis}
  \bigl\langle\xi^{m,m},P_X(Y)\bigr\rangle\varepsilon_Y
  =f^{*}(\xi^{m,m})
  \qquad
  \bigl(\xi^{m,m}\in H^{m}(X,\Omega_X^{m})\bigr),
\end{equation}
\footnote{印刷文本写作 $f_{*}(\xi^{m,m})$。由于
$\xi^{m,m}\in H^{m}(X,\Omega_X^{m})$，且等式两端均属于
$H^{m}(Y,\Omega_Y^{m})$，这里需要的映射是 $f^{*}$。}
其中 $\varepsilon_Y$ 是 $H^{m}(Y,\Omega_Y^{m})$ 中的基本元；对于非奇异
射影代数簇，该式给出了 $Y$ 所对应上同调类的一个新定义。为证明定理 4，
借助定理 3，把同态
\[
  H^{m}(X,\Omega_X^{m})
  \longrightarrow H^{m}(Y,\Omega_Y^{m})
  =H^{m}(X,\Omega_Y^{m})
\]
的转置解释为同态
\begin{equation}\tag{7.5}
\begin{split}
  H^{n-m}(X,\Omega_X^{n-m})
  &\simeq
  \operatorname{Ext}_{\Osh_X}^{n-m}(X;\Omega_X^{m},\Omega_X^{n})\\
  &\longleftarrow
  \operatorname{Ext}_{\Osh_X}^{n-m}(X;\Omega_Y^{m},\Omega_X^{n})
  \simeq
  \operatorname{Hom}_{\Osh_X}(X;\Omega_Y^{m},\Omega_Y^{m}).
\end{split}
\end{equation}
可以验证，$H^{m}(Y,\Omega_Y^{m})$ 的对偶中的元素 $1_Y$，与右端中对应于
$\Omega_Y^{m}$ 恒等自同态的元素等同；另一方面，该元素在
$H^{n-m}(X,\Omega_X^{n-m})$ 中的像正是 $P_X(Y)$。这些相容性本可在第 4 节
就给出；事实上，它们可以对任意非奇异代数簇表述，并且确实成立。例如，第二个
相容性来自下列典范自同态图的交换性：

\src{189}
\begin{equation}\tag{7.6}
\resizebox{\textwidth}{!}{$
\left\{
\begin{array}{ccccc}
\operatorname{Ext}_{\Osh_X}^{n-m}(X;\Omega_X^{m},\Omega_X^{n})
&\longleftarrow&
\operatorname{Ext}_{\Osh_X}^{n-m}(X;\Omega_Y^{m},\Omega_X^{n})
&\simeq&
\operatorname{Hom}_{\Osh_X}(X;\Omega_Y^{m},\Omega_Y^{m})\\[3pt]
\Big\downarrow&&&&\Big\downarrow\\[3pt]
\operatorname{Ext}_{\Osh_X}^{n-m}(X;\Osh_X,\Omega_X^{n-m})
&\longleftarrow&
\operatorname{Ext}_{\Osh_X}^{n-m}(X;\Osh_Y,\Omega_X^{n-m})
&\simeq&
\operatorname{Hom}_{\Osh_X}(X;\Omega_X^{m},\Omega_Y^{m})
\end{array}
\right.
$}
\end{equation}

由此得到紧致定向流形上庞加莱对偶形式体系的一个完全对应物。特别地，定理 4
可以确定 $X\times X$ 的对角线所对应的上同调类。再由一个熟知的论证，例如可得
如下 Lefschetz 公式。

\result{定理}{5} 设 $f$ 是非奇异射影代数簇 $X$ 的一个自同态，并且 $f$ 的
所有不动点的重数均为 1。则这些不动点的数目模 $k$ 的特征，同余于 $f$ 在各
$H^{i}(X)$ 上所诱导自同态之迹的交错和。

对 $f$ 所作的限制源于第 4 节注记中提到的困难。不过应当注意，若 $f$ 与一个
所有不动点重数均为 1 的自同态\fquote{同伦}，Lefschetz 公式仍然成立。

\fsection{8}{对偶定理的推广}

设 $X$ 是一个非奇异代数簇，并且 $X$ 上每个凝聚代数层 $F$ 都同构于某个局部
自由凝聚代数层的商（当 $X$ 是射影空间中的局部闭子集时即属此情形）。于是
$X$ 上每个凝聚代数层 $F$ 都具有一个由局部自由层组成的有限消解 $L$；对于
任意两个这样的消解，总能找到第三个消解以及从它到前两个消解、且与增广相容的
同态。同样，若 $L$ 是 $F$ 的局部自由有限消解，并给定同态
$F'\longrightarrow F$，则存在 $F'$ 的局部自由有限消解 $L'$ 以及与
$F'\longrightarrow F$ 相容的同态 $L'\longrightarrow L$；若
$F'\longrightarrow F$ 是满射，还可以令后一同态为满射。由此，对于给定的两个
整数 $r,s\geq0$，可以定义两个关于诸变元
$A_1,\ldots,A_r$；$B_1,\ldots,B_s$ 的上同调多重函子；

\src{190}
这些变元在 $X$ 上凝聚代数层范畴中变化，而两个函子分别取值于 $X$ 上凝聚
代数层范畴和 $H^{0}(X,\Osh_X)$-模范畴；其定义式为
\begin{equation}\tag{8.1}
\resizebox{\textwidth}{!}{$
\left\{
\begin{aligned}
\underline{T}_{r}^{s*}(A_1,\ldots,A_r;B_1,\ldots,B_s)
  &={}H^{*}\!\left(
    \Homsh_{\Osh_X}\!\left(
      L(A_1)\otimes\cdots\otimes L(A_r),
      L(B_1)\otimes\cdots\otimes L(B_s)
    \right)\right),\\[3pt]
T_{r}^{s*}(A_1,\ldots,A_r;B_1,\ldots,B_s)
  &={}\mathbf R^{*}\Gamma\!\left(
    \Homsh_{\Osh_X}\!\left(
      L(A_1)\otimes\cdots\otimes L(A_r),
      L(B_1)\otimes\cdots\otimes L(B_s)
    \right)\right).
\end{aligned}
\right.
$}
\end{equation}
在该公式中，$L(F)$ 表示凝聚代数层 $F$ 的一个局部自由有限消解，
$\mathbf R^{*}\Gamma(K)$ 表示空间 $X$ 关于层复形 $K$ 的超上同调。若 $r$ 或
$s$ 为零，就分别用 $\Osh_X$ 代替诸 $L(A_i)$ 或诸 $L(B_j)$ 的张量积。特别地，
$\underline{T}_{0}^{0}$ 与 $T_{0}^{0}$ 是零变元的分次函子：
$\underline{T}_{0}^{0}$ 只在次数 0 非零，此时它是层 $\Osh_X$；而
$T_{0}^{0}$ 等同于 $H^{*}(X,\Osh_X)$。(8.1) 右端不依赖所选消解：对第一行
这无论如何是显然的（因为问题是局部的）；对第二行，则由此前的一般性说明以及
层复形
\[
  K=\Homsh_{\Osh_X}\bigl(L(A_1)\otimes\cdots,
      L(B_1)\otimes\cdots\bigr)
\]
的超上同调谱序列推出。该谱序列收敛于 $X$ 关于 $K$ 的超上同调，其初始项是
$H^{p}(X,H^{q}(K))$，即
\begin{equation}\tag{8.2}
  E_{2}^{p,q}
  =H^{p}\!\left(X,
    (\underline{T}_{r}^{s})^{(q)}
      (A_1,\ldots,A_r;B_1,\ldots,B_s)\right).
\end{equation}
由此可见，该谱序列本身也不依赖所选消解，而其收敛项就是 (8.1) 中相应的
左端。又因每个正合列
$0\longrightarrow F'\longrightarrow F\longrightarrow F''\longrightarrow0$
都能由局部自由有限复形的一个正合列来消解，故容易定义关于各变元
$A_i,B_j$ 的连接算子。

在函子系 $\underline{T}_{r}^{s*}$ 以及 $T_{r}^{s*}$ 中，可以定义类似于
张量演算的运算；它们由定义式 (8.1) 立即给出。例如，有一个复合运算
（它推广了第 2 节所考察的复合）
\begin{equation}\tag{8.3}
\resizebox{\textwidth}{!}{$
\begin{aligned}
&T_{r}^{s*}(A_1,\ldots,A_r;B_1,\ldots,B_s)
\times T_{r'}^{s'*}(A'_1,\ldots,A'_{r'};B'_1,\ldots,B'_{s'})\\
&\qquad\longrightarrow
T_{r+r'}^{s+s'*}
(A_1,\ldots,A_r,A'_1,\ldots,A'_{r'};
 B_1,\ldots,B_s,B'_1,\ldots,B'_{s'}),
\end{aligned}
$}
\end{equation}
它满足显然的结合律，并与函子性同态、连接同态和谱序列相容。同样，还有一些
对称运算，其具体写法留给读者。再者，每当某个变元 $A_i$

\src{191}
等于某个变元 $B_j$ 时，就有一个缩并运算：
\begin{equation}\tag{8.4}
\resizebox{\textwidth}{!}{$
\begin{aligned}
&T_{r}^{s*}(A_1,\ldots,A_{i-1},C,A_{i+1},\ldots,A_r;
 B_1,\ldots,B_{j-1},C,B_{j+1},\ldots,B_s)\\
&\qquad\longrightarrow
T_{r-1}^{s-1*}(A_1,\ldots,\widehat{A_i},\ldots,A_r;
 B_1,\ldots,\widehat{B_j},\ldots,B_s).
\end{aligned}
$}
\end{equation}
此外，若变元 $A_i$ 是局部自由层，则可以消去它，只需把另一个 $A_j$
（$j\neq i$）换成 $A_j\otimes A_i$，或把某个 $B_k$ 换成
$B_k\otimes A_i'$（其中 $A_i'=\Homsh_{\Osh_X}(A_i,\Osh_X)$）；当某个变元
$B_i$ 局部自由时，也有类似的计算规则。特别地，总可以消去等于 $\Osh_X$ 的
变元。若除至多一个 $B_i$ 外所有变元都局部自由，刚才的规则给出函子性同构
\begin{equation}\tag{8.5}
  T_{r}^{s*}(A_1,\ldots,A_r;B_1,\ldots,B_s)
  =H^{*}\bigl(X,A_1'\otimes\cdots\otimes A_r'
    \otimes B_1\otimes\cdots\otimes B_s\bigr)
\end{equation}
（因为这可归结到 $r=0,s=1$ 的情形，而该情形是立即可见的；也可以直接使用
初始项为 (8.2) 的谱序列）。对 $\underline{T}_{r}^{s}$ 也可定义与前述运算
相应的运算。如此引入的各运算之间的关系，与张量演算中类似运算之间的关系相同。

设 $n$ 是 $X$ 的维数。依次应用张量复合 (8.3)，再对重复出现的变元施行缩并
(8.4)，便得到配对
\begin{equation}\tag{8.6}
\resizebox{\textwidth}{!}{$
  (T_{r}^{s})^{p}(A_1,\ldots;B_1,\ldots)
  \times
  (T_{s}^{r})^{n-p}(B_1,\ldots;A_1,\ldots,A_r\otimes\Omega_X^{n})
  \longrightarrow H^{n}(X,\Omega_X^{n}).
$}
\end{equation}

\result{定理}{6} 若 $X$ 是非奇异射影代数簇，则配对 (8.6) 都是对偶。

这完全形式地由定理 3 的推论得出。事实上，从该推论容易推出：若 $K$ 是由
局部自由凝聚代数层组成的复形，则 $X$ 关于 $K$ 的超上同调，与 $X$ 关于
$K'\otimes\Omega_X^{n}$ 的超上同调，通过自然配对
\begin{equation}\tag{8.7}
  \mathbf R^{p}\Gamma(K)
  \times\mathbf R^{n-p}\Gamma(K'\otimes\Omega_X^{n})
  \longrightarrow\mathbf R^{n}\Gamma(\Omega_X^{n})
  =H^{n}(X,\Omega_X^{n}).
\end{equation}
互为对偶。这可由初始项为
$H^{p}(X,H^{q}(K))$\footnote{印刷文本写作
$H^{p}(H^{q}(X,K))$；初始项应为谱序列 (8.2) 的初始项，即
$H^{p}(X,H^{q}(K))$。}
的谱序列，以及 $K'\otimes\Omega_X^{n}$ 的类似谱序列看出。利用定义 (8.1)，
定理 6 即由上述结果推出。

\src{192}
\subsection*{注记}

\noindent 1\textsuperscript{o} 在定理 6 之前的定义中，并不需要假设 $X$
非奇异，因为只使用有限消解并非不可缺少。但是，当 $X$ 奇异时，不能再先验地
断言
$(\underline{T}_{r}^{s})^{p}(A_1,\ldots;B_1,\ldots)$ 是凝聚层，因为在层复形
\[
  \Homsh_{\Osh_X}\bigl(L(A_1)\otimes\cdots,L(B_1)\otimes\cdots\bigr)
\]
中，每个给定的总次数都会有无穷多个分量。

\medskip
\noindent 2\textsuperscript{o} 容易验证，在公式 (8.1) 中，可以把某个
$L(B_i)$ 换成 $B_i$。结合命题 3，这就表明
\begin{equation}\tag{8.8}
\left\{
\begin{aligned}
  \underline{T}_{1}^{1*}(A;B)&=\Extsh_{\Osh_X}^{*}(A,B),\\
  T_{1}^{1*}(A;B)&=\operatorname{Ext}_{\Osh_X}^{*}(X;A,B).
\end{aligned}
\right.
\end{equation}
特别地，在 (8.6) 中取 $r=s=1$、$A_1=\Osh_X$，便重新得到定理 3。公式
(8.8) 还蕴含
\[
  T_{0}^{1*}(B)=H^{*}(X,B),
  \qquad
  T_{1}^{0*}(A)=\operatorname{Ext}_{\Osh_X}^{*}(X;A,\Osh_X).
\]

\medskip
\noindent 3\textsuperscript{o} 由公式 (8.1) 可见，一般而言，函子
$\underline{T}_{r}^{s*}$ 和 $T_{r}^{s*}$ 同时具有正次数和负次数的分量。
利用注记 2 可见：若 $X$ 的维数为 $n$，则当 $s>0$ 时，
$\underline{T}_{r}^{s*}$ 的非零分量介于 $-(s-1)n$ 与 $rn$ 之间；当
$s=0$ 时，介于 $0$ 与 $rn$ 之间。对于 $T_{r}^{s*}$，当 $s>0$ 时，其
非零分量介于 $-(s-1)n$ 与 $(r+1)n$ 之间；当 $s=0$ 时，介于 $0$ 与
$(r+1)n$ 之间（而且，若无误，当 $r>0$ 时，甚至分别介于 $-(s-1)n$
与 $rn$ 之间以及 $0$ 与 $rn$ 之间）。

\section*{参考文献}

\noindent[1] Cartier (Pierre). --- Les groupes
$\operatorname{Ext}^{s}(A,B)$, Séminaire A. Grothendieck : Algèbre homologique,
t. 1, 1957, n\textsuperscript{o} 3.

\noindent[2] Grothendieck (Alexandre). --- Sur quelques points d'algèbre
homologique, \emph{Tohoku Math. J.}, t. 9, 1957, p. 119--183.

\noindent[3] Serre (Jean-Pierre). --- Faisceaux algébriques cohérents,
\emph{Annals of Math.}, t. 61, 1955, p. 197--278.

\noindent[4] Serre (Jean-Pierre). --- Sur la dimension homologique des anneaux
et des modules noethériens, \emph{Proc. Intern. Symp. on Alg. Number Theory}
[1955, Tokyo et Nikko]. --- Tokyo, Science Council of Japan, 1956 ;
p. 175--189.

\clearpage
\src{193}
\section*{补记}

第 13 页注记中指出的困难目前已经完全解决，这是由于对偶定理已推广到任意
代数簇（允许任意奇点）。关于这一理论——它可以在\fquote{概形理论}的一般框架
中表述——我们请读者参阅 J. Dieudonné 与 A. Grothendieck 当时正在编写的
《Éléments de Géométrie algébrique》。下列文献给出了一些说明：

\begin{quote}
Grothendieck (A.). --- Cohomology theory of algebraic varieties,
\emph{Congrès int. des Mathématiciens}, 1958, Edinburgh (à paraître).
\end{quote}

\begin{flushright}
[1959年4月]
\end{flushright}

\begin{center}\rule{36mm}{0.5pt}\end{center}

\par
